%%%%%%%%%%%%%%%%%%%%%%%%%%%%%%%%%
% Conclusion                    %
%%%%%%%%%%%%%%%%%%%%%%%%%%%%%%%%%

\section{Conclusions and Future Work}

This chapter closes the report. It first draws the conclusions of the work against the objectives and the research questions stated in the introduction, and then outlines the directions in which the framework and the delivered platform could be extended.

\subsection{Conclusions}

This project set out to address a concrete and pressing problem: the Tunisian banking sector, under the dual pressure of the BCT Circulaire N\textdegree{}2025-08 and the international BCBS 239 standard, had no standardized, rigorous, and regulation aligned means of measuring and steering its data maturity. Existing maturity models were either too generic to capture the specific obligations of a supervised bank, or too abstract to be applied without extensive interpretation, and none combined a defensible measurement method with a tool that management could operate directly. The work delivered a complete answer to that problem through two connected artifacts, a hybrid data maturity framework and the DataPilot tool that operationalizes it, both designed and validated through a Design Science Research methodology combined with an Agile, iteration based development process.

The study met each of its five research objectives, and it is useful to state precisely how. The first objective, to diagnose the current state of the sector, was met through a diagnostic survey of forty four banking professionals, which established that data maturity is concentrated at the emerging level, that is Level 2 on the five level scale, and that the most significant gaps lie in the Governance and Data Quality dimensions. The second objective, to select an appropriate set of reference models, was met by analyzing fifteen candidate maturity frameworks against a common set of selection criteria and retaining the five most relevant, namely DAMA-DMBOK, DCAM, CMMI-DMM, TDWI, and Gartner, with an explicit record of what was retained and what was excluded from each. The third objective, to design a framework adapted to the sector, produced a hybrid structure of five dimensions, twelve sub-dimensions, and forty seven evidence based indicators, with dimension weights of twenty five, twenty, twenty, twenty, and fifteen percent for Governance, Data Quality, Architecture and Access, Analytics and Tools, and Skills and Culture respectively, and with thirteen non-skippable regulatory indicators anchored in the BCT Circulaire N\textdegree{}2025-08 and BCBS 239 requirements. The fourth objective, to make the measurement auditable, was met by an aggregation method that averages indicator scores into sub-dimension scores, sub-dimension scores into dimension scores, and dimension scores into a Global Maturity Index on a one to five scale, governed by an evidence cap that limits any unsubstantiated high score to two, by a bounded skip logic that permits at most a fifth of a dimension's indicators to be deferred, and by a rule that forbids skipping any of the thirteen regulatory indicators. The fifth objective, to deliver a usable instrument, was met by DataPilot itself, which delivers the six screen assessment workflow that automates the scoring and presents the maturity profile, the gap analysis, the three phase improvement roadmap, and the regulatory compliance view, and which, in the final iteration, became a multi-tenant platform with four roles, EY Admin (owner), Super Admin, Admin, and Analyst, through which EY administers the instrument across banks and each bank conducts its own solo and group assessments. One clarification is owed on the fourth operational problem stated in Chapter 1, which called for a tool that also simulates the business benefits of a data transformation: this simulation need is addressed today through the setting of a target maturity level, the visualization of the gap that separates the current profile from that target, and the deterministic three phase roadmap that sequences its closure, while the quantitative what-if simulation of business gains is the identified extension developed in the future work below, an extension that the deterministic aggregation chain already makes straightforward to add.

The three research questions were answered correspondingly. The current state of data maturity in Tunisian banks was characterized, with its gaps relative to international standards and regulatory requirements made measurable and explicit rather than described in general terms. A framework adapted to the sector was designed by combining the five benchmark frameworks with the survey findings, with each reference model contributing a distinct and non-overlapping role, from the knowledge area structure of DAMA-DMBOK to the maturity ladder of CMMI-DMM and the transformation lens of Gartner. And the capacity of an interactive steering tool to support strategic decision making was demonstrated through DataPilot, which converts a static assessment into an actionable instrument by automating the scoring, renormalizing the weights over the dimensions actually answered, prioritizing the gaps by severity, and exposing the regulatory readiness of an institution through a dedicated compliance view. The demonstrated capacity now extends to collaborative assessment at the scale of an institution, with the responsibility for each dimension distributed to the department that holds its evidence and a coordinator consolidating the contributions into one auditable submission.

Beyond meeting its objectives, the project delivered three contributions of lasting value. The first is a conceptual contribution, a banking specific and regulation aligned hybrid framework that unifies five established reference models into a single coherent structure and binds it to the Tunisian supervisory context. The second is a methodological contribution, a rigorous and auditable evidence based scoring method in which the evidence cap, the skip limit, and the regulatory non-skippable rule together make a maturity score defensible before an auditor rather than merely indicative. The third is a practical contribution, an operational tool that bank management can use directly, without specialist intermediation, to obtain a maturity profile, a prioritized roadmap, and a compliance position in a single session, delivered as a multi-tenant platform in which EY maintains the master reference instrument, each bank tailors its own copy, and the four roles collaborate under database-enforced tenant isolation. The worked example developed in this report, in which a representative institution scores a Global Maturity Index of 2.19 out of 5, equivalent to 44 percent and to the emerging level, illustrates how the three contributions combine to turn a set of survey responses into an interpretable and actionable diagnosis. Together, these contributions provide the Tunisian banking sector with a rigorous and actionable instrument to assess its data maturity and to steer its digital transformation in alignment with regulatory expectations.

Three external constraints shaped the work, and in each case for the better. First, no direct instrument existed for measuring skills and culture, which forced the proxy only rule for the fifth dimension; the consequence is that the least observable dimension is explicitly labeled as proxy based and deliberately bounded at fifteen percent of the index, instead of being silently unreliable. Second, the confidentiality expectations attached to banking data first forced a fully client-side architecture, in which every response, score, and evidence reference remained in the browser local storage of the assessing institution; when the platform of the final iteration made a shared backend unavoidable, the same expectations forced the isolation discipline that replaced local custody: every row is confined to its own bank by row-level security, access is invite-only with no public sign-up, no user can escalate their own role, and submissions are immutable once recorded. The custody argument became a tenant isolation argument, and the constraint sharpened the design a second time. Third, the absence of access to production bank data confined the evaluation to simulated and anonymized scenarios, which concentrated it on the determinism and fidelity of the instrument, precisely the properties that a future field pilot will inherit rather than re-establish. These constraints delimit the claims the report can make, but they also sharpened the design, and their honest documentation is itself part of the contribution. The recurring lesson of the work, established in the discussion and worth restating here, is that the score is trustworthy precisely where the method constrains it: the evidence cap, the skip ceilings, and the non-skippable regulatory rule are not restrictions on the instrument but the source of its defensibility.

\subsection{Future Work}

While the framework and the tool are complete and functional, several directions would extend their value and address the limitations acknowledged in this report. The directions below are ordered roughly by immediacy, from the empirical validation that would consolidate the present results to the broader extensions that would enlarge the scope of the work. None of them require reworking the core artifacts; each builds on the structural model, the scoring methodology, and the tool as they stand.

\subsubsection{Validation on Real Institutional Data}

The tool was evaluated on simulated and anonymized banking data rather than on a full production deployment inside a supervised institution. The most immediate priority is therefore to conduct pilot assessments with one or more real banks, ideally spanning institutions of different sizes and maturity levels so that the framework is exercised across the full range of the one to five scale. Such pilots would serve three purposes at once. They would calibrate the forty seven indicators against field reality, confirming that each rubric level is distinguishable and that the evidence requirements are neither too lax nor too demanding for the kind of documentation banks actually hold. They would confirm the usability of the questionnaire for bank staff who are not data specialists, in particular the clarity of the hints and the workability of the evidence cap when real evidence must be attached. And they would validate the relevance of the improvement roadmaps by checking that the actions surfaced for a low scoring dimension correspond to the steps practitioners would themselves prioritize. A structured validation workshop with governance officers, data owners, and compliance staff would then refine the rubrics, the weighting, and the mapping of indicators to the BCT and BCBS 239 requirements, closing the loop between the designed artifact and its intended field of use.

\subsubsection{Strengthening the Skills and Culture Dimension}

The fifth dimension, Skills and Culture, is currently assessed through proxy indicators only, since no direct self assessment instrument was available during the project and since cultural attributes are intrinsically harder to observe than the documentary artifacts that anchor the other four dimensions. As a consequence, this dimension carries none of the thirteen regulatory indicators and its eight indicators infer culture from measurable surrogates such as training budgets, the ratio of data roles to headcount, recent certifications, and the existence of communities of practice. Future work could develop a dedicated instrument, for example a structured staff survey or a data literacy assessment administered across business units, that would measure the culture directly rather than approximating it through investment proxies. Combining such an instrument with the existing proxy indicators would improve the construct validity of the dimension, allow the proxy and direct measures to be cross checked against each other, and give the fifteen percent weight assigned to Skills and Culture a firmer empirical basis.

\subsubsection{Sector Level Benchmarking and Anonymized Comparison}

Because the framework produces a Global Maturity Index that is comparable across institutions by construction, using the same dimensions, indicators, weights, and aggregation rules for every bank, a natural extension is a sector benchmarking capability that would allow each institution to position itself against an anonymized sector distribution rather than against an absolute scale alone. A bank could then see not only that it sits at the emerging level, but where it stands relative to its peers on each dimension, which sharpens the interpretation of a score and helps management judge the urgency of a gap. The platform delivered in the final iteration now provides the collection mechanism this capability presupposed, since every finalized assessment is recorded as an immutable submission in a central store that EY can review across banks; what remains to build is the aggregation layer above it. Realizing this capability would still require a governance arrangement for the aggregation of results, with careful attention to anonymization so that no individual institution can be identified from the aggregates, ideally established in coordination with the regulator given its natural role as a trusted aggregator. Such an arrangement would transform the tool from an institutional instrument into a sector observatory, and would give the supervisor a consolidated, evidence based view of data maturity across the banks it oversees.

\subsubsection{Scenario Simulation and Target Setting}

DataPilot already supports the selection of a target maturity level against which the current profile is compared. A valuable extension would be a scenario simulation capability that estimates the effect of specific remediation actions on the Global Maturity Index before any of them is undertaken. Because the aggregation chain is fully deterministic, from indicator to sub-dimension to dimension to the weighted and renormalized index, the tool already contains everything needed to recompute a projected score when selected indicators are raised to a higher rubric level. Exposing this as an interactive what if analysis would let management compare investment options by their projected impact on maturity and, just as importantly, on regulatory compliance, since the same simulation could show how many of the thirteen regulatory indicators would move from non compliant to compliant. Layered over successive assessments, this would let an institution plan a maturity trajectory over time and choose the sequence of actions that closes its regulatory gaps most efficiently.

\subsubsection{AI-Assisted Roadmap Generation}

The improvement roadmap is currently generated deterministically, by banding each dimension score and surfacing a curated set of recommended actions, which keeps every suggestion reproducible and auditable. A natural extension is to let a language model rephrase or enrich these actions for the specific institution being assessed. This was deliberately left as future work rather than delivered, because sending a bank's assessment content to an external model raises data-residency and confidentiality questions that must be resolved first, for example by using a model hosted within the bank's own jurisdiction or by restricting the model's input to non-sensitive dimension labels. Any such feature would also have to preserve the determinism and auditability of the underlying roadmap, so that the generated text remains an optional presentation layer over a reproducible core rather than the source of the recommendation itself.

\subsubsection{From Delivered Platform to Enterprise Integration}

The backend, the persistence of multiple assessments over time, and the multi-user collaboration that earlier stages of this project could only list as future work were delivered by the fourth iteration, described in Chapter 4, and are therefore recorded here as achievements rather than prospects. The platform stores every finalized result as an immutable submission in a central database, so that an institution's assessments accumulate over time and successive maturity positions can be compared; the group assessment lets each department answer the indicators of its own dimensions within one shared evaluation, which is precisely the collaboration in which data owners answer their own domain; and the authorship of every group answer is recorded, giving the audit trail its first concrete form. What remains is the enterprise integration that would embed the platform in the daily infrastructure of a bank and of the advisory practice, and four directions stand out. First, single sign-on and enterprise identity federation would connect the invite based accounts to the directories that banks already operate, so that joiners, movers, and leavers are governed by the institution's existing identity lifecycle rather than by manual invitations alone. Second, audit dashboards would surface the authorship and history that the platform already records as a reviewable trail, letting an internal auditor see who answered each indicator of a group assessment, when a draft was finalized, and how the submissions of an institution have evolved between assessments. Third, sector benchmarking, discussed above, can now be computed over the submissions that accumulate in the platform, since the finalized results of every bank are already collected under one governed store; the remaining work is the anonymized cross bank aggregation and the governance arrangement around it. Fourth, automated evidence checking would assist reviewers by flagging evidence texts that are empty, duplicated across indicators, or inconsistent with the claimed rubric level, extending to the review stage the discipline that the evidence cap already imposes at the moment of scoring.

\subsubsection{Extension to Adjacent Sectors}

The framework is deliberately specific to the banking sector and to the Tunisian regulatory environment, and this specificity is a source of its rigor. Its architecture, however, separates the structural model, that is the five dimensions, the twelve sub-dimensions, and the general purpose indicators, from the regulatory anchoring carried by the thirteen indicators tied to the BCT Circulaire N\textdegree{}2025-08 and BCBS 239. This separation is deliberate and it makes the framework portable. Future work could adapt it to adjacent regulated sectors such as insurance by re-anchoring the regulatory indicators to the relevant supervisory texts, for example the obligations of the insurance regulator, while retaining the structural backbone, the scoring methodology with its evidence cap and skip logic, and the tool itself. The same approach would extend the framework to other jurisdictions by substituting the local banking regulation for the Tunisian circular. In each case the engineering effort would concentrate on the regulatory indicators and their references, leaving the bulk of the artifact, and the accumulated design rationale behind it, intact.

\subsection{Closing Remark}

This project demonstrates that the gap between a rigorous academic maturity model and a usable management instrument can be closed. By grounding the framework in established research, anchoring it in binding regulation, and delivering it through an operational tool and a governed platform through which an advisor can operate it across the sector, the work provides the Tunisian banking sector with both a measure of where it stands and a path toward where regulation and competition require it to be.
